
\setcounter{equation}{0}
\setcounter{Q}{0}

\begin{center}
\begin{minipage}{0.48\textwidth}
\includegraphics[width=\textwidth]{/home/jeff/Documents/IUTOrsay/1-Cours/4p13-Alliages/Cours-TD/Exos/Binaires-Lecture/Homogene-Ex03/figures/Fig01}
\end{minipage}\hfill\begin{minipage}{0.48\textwidth}
\captionof{figure}{Diagramme binaire isobare, à une pression de \SI{1}{\bar}, du mélange \ce{Cu}-\ce{Ni}.}
[La composition est donnée par la fraction massique en nickel dans le mélange]
\end{minipage}
\end{center}

Le diagramme binaire isobare solide-liquide du système cuivre-nickel est donné à la figure ci-dessus.

\Q Déterminer les températures de fusion de ces deux métaux.

\Q Identifier les deux courbes, les trois domaines qu'elles définissent et la variance particulière pour chacun d'eux.

\Q Un mélange liquide de cuivre et de nickel commence à se solidifier à \SI{1200}{\celsius}. En déduire : 
\begin{enumerate}
\item la composition du mélange et celle du premier cristal qui apparaît ;
\item la composition de la dernière goutte qui disparaît.
\end{enumerate}

\Q Un mélange liquide de cuivre et de nickel a une fraction molaire en cuivre $x^l (\ce{Cu}) = 0.40$. Déterminer : 
\begin{enumerate}
\item la composition du mélange en fraction massique ;
\item la température de début de solidification du mélange ; 
\item la composition, en fraction massique, de la dernière goutte qui disparaît.
\end{enumerate}

\paragraph*{Données} 
\begin{align*}
M(\ce{Ni}) &= 58.7 ~ \si{\gram\per\mole} \\
M(\ce{Cu}) &= 63.5 ~ \si{\gram\per\mole} \\
\end{align*}